
\usepackage[ansinew]{inputenc}


\usepackage{doi}

% Number only equations which are referenced in the text
% The command \eqref needs to be "robustified" in order for
% the numbering to work OK when \eqref is used in moving
% places, such as figure captions. As an alternative, you
% may also prepend eqref with a \protect command, as follows:
% \caption{ ... \protect\eqref{eq:tag} ... }
% No {} around the \eqref command though.
% I actually use the alternative, because the \MakeRobust
% command doens't seem to work.
\usepackage{amsmath}
%\usepackage{fixltx2e,amsmath}
%\MakeRobust{\eqref}
%\usepackage{mathtools}
%\mathtoolsset{showonlyrefs}

% For the jury : number lines -- doesn't work very well though ...
%\usepackage{xcolor}
%\usepackage[displaymath]{lineno}
%% Running line numbers:
%% Same, but that reset on every page:
%\pagewiselinenumbers
%\switchlinenumbers
%% Number only every 5:th line:
%\modulolinenumbers[5]
%\setlength\linenumbersep{5pt}
%\renewcommand\linenumberfont{\normalfont\tiny\sffamily\color{gray}}

\usepackage{mathtools}
\DeclarePairedDelimiter{\ceil}{\lceil}{\rceil}
\DeclarePairedDelimiter{\floorr}{\lfloor}{\rfloor}


\usepackage{amsfonts}
\usepackage{amsthm}
\usepackage{amssymb}

\usepackage{color}
\usepackage{url}
\usepackage{bm} % bold math fonts with \bm{math expr}
\usepackage{dsfont} % blackboard bold ones

\usepackage{graphicx}


%\usepackage[lmargin=3cm,rmargin=3cm,bottom=3cm,top=3cm]{geometry}


% For fancy tables
% https://tex.stackexchange.com/questions/72945/how-to-merge-cells-vertically
\usepackage{adjustbox}
\usepackage{multirow}
\usepackage{longtable}



\usepackage{tikz}
\definecolor{mycolor1}{rgb}{0.105882,0.619608,0.466667}
\definecolor{mycolor2}{rgb}{0.85098,0.372549,0.00784314}
\definecolor{mycolor3}{rgb}{0.458824,0.439216,0.701961}
\definecolor{mycolor4}{rgb}{0.905882,0.160784,0.541176}
\definecolor{mycolor5}{rgb}{0.4,0.65098,0.117647}
\definecolor{mycolor6}{rgb}{0.65098,0.462745,0.113725}
\definecolor{mycolor7}{rgb}{0.901961,0.670588,0.00784314}
\definecolor{mycolor8}{rgb}{0.4,0.4,0.4}
\definecolor{mycolor9}{rgb}{0.301961,0,0.294118}
\definecolor{mycolor10}{rgb}{0.0313725,0.25098,0.505882}


\usetikzlibrary{calc}        

\makeatletter
\newif\ifmygrid@coordinates
\tikzset{/mygrid/step line/.style={line width=0.80pt,draw=gray!80},
         /mygrid/steplet line/.style={line width=0.25pt,draw=gray!80}}
\pgfkeys{/mygrid/.cd,
         step/.store in=\mygrid@step,
         steplet/.store in=\mygrid@steplet,
         coordinates/.is if=mygrid@coordinates}
\def\mygrid@def@coordinates(#1,#2)(#3,#4){%
    \def\mygrid@xlo{#1}%
    \def\mygrid@xhi{#3}%
    \def\mygrid@ylo{#2}%
    \def\mygrid@yhi{#4}%
}
\newcommand\DrawGrid[3][]{%
    \pgfkeys{/mygrid/.cd,coordinates=true,step=1,steplet=0.2,#1}%
    \draw[/mygrid/steplet line] #2 grid[step=\mygrid@steplet] #3;
    \draw[/mygrid/step line] #2 grid[step=\mygrid@step] #3;
    \mygrid@def@coordinates#2#3%
    \ifmygrid@coordinates%
        \draw[/mygrid/step line]
        \foreach \xpos in {\mygrid@xlo,...,\mygrid@xhi} {%
          (\xpos,\mygrid@ylo) -- ++(0,-3pt)
                              node[anchor=north] {$\xpos$}
        }
        \foreach \ypos in {\mygrid@ylo,...,\mygrid@yhi} {%
          (\mygrid@xlo,\ypos) -- ++(-3pt,0)
                              node[anchor=east] {$\ypos$}
        };
    \fi%
}
\makeatother




% use in a table cell to authorize linebreaks with \\ %
% http://tex.stackexchange.com/questions/2441/how-to-add-a-forced-line-break-inside-a-table-cell
\newcommand{\specialcell}[2][c]{%
  \begin{tabular}[#1]{@{}#1@{}}#2\end{tabular}}

% for sidewaystable : tables in landscape mode
%\usepackage{rotating}

\usepackage{psfrag}

\usepackage{algorithm}
\usepackage{algpseudocode}
% Do-While construct
\algdef{SE}[DOWHILE]{Do}{doWhile}{\algorithmicdo}[1]{\algorithmicwhile\ #1}%

% section numbering in margin
%\usepackage{sectsty}
%\makeatletter\def\@seccntformat#1{\protect\makebox[0pt][r]{\csname the#1\endcsname\hspace{12pt}}}\makeatother

\usepackage{mathrsfs} % for \mathscr{ABC} style (curly calligraphy)


%\usepackage[colorlinks=true,bookmarks=false]{hyperref}

%\usepackage[round,compress]{natbib}

%\usepackage[square,numbers,compress]{natbib}

%%%%% Activate these in track change mode:
%\newcommand{\add}[1]{{\color{blue}{#1}}}
%\newcommand{\remove}[1]{{\footnote{\color{red}{[removed: #1]}}}}
%\newcommand{\removesafe}[1]{{\scriptsize{\color{red}{[removed: #1]}}}}
%\newcommand{\change}[2]{\remove{#1}\add{#2}}
%\newcommand{\changesafe}[2]{\removesafe{#1}\add{#2}}
%%%%% Activate these in release mode:
\newcommand{\add}[1]{{#1}}
\newcommand{\remove}[1]{}
\newcommand{\removesafe}[1]{}
\newcommand{\change}[2]{{#2}}
\newcommand{\changesafe}[2]{{#2}}

\usepackage{subfigure}


% show labels in the margin (for debugging)
% \usepackage[right]{showlabels}



\newcommand{\ER}{{Erd\H{o}s-R\'enyi }}
\newcommand{\argmin}[1]{\underset{#1}{\operatorname{argmin}}}
\newcommand{\argmax}[1]{\underset{#1}{\operatorname{argmax}}}
\newcommand{\transpose}{^\top\! }
\newcommand{\eig}{{\mathrm{eig}}}
\newcommand{\mle}{{\mathrm{MLE}}}
\newcommand{\mlez}{{\mathrm{MLE0}}}
\newcommand{\inner}[2]{\left\langle{#1},{#2}\right\rangle}
\newcommand{\innersmall}[2]{\langle{#1},{#2}\rangle}
\newcommand{\innerbig}[2]{\big\langle{#1},{#2}\big\rangle}
\newcommand{\MSE}{\mathrm{MSE}}
\newcommand{\Egood}{E_\textrm{good}}
\newcommand{\Ebad}{E_\textrm{bad}}
\newcommand{\trace}{\mathrm{Tr}}
\newcommand{\Trace}{\mathrm{Tr}}
\newcommand{\spann}{\mathrm{span}}
\newcommand{\im}{\mathrm{im}}
%\newcommand{\skeww}[1]{\left\lceil #1 \right\rfloor}
\newcommand{\skeww}[1]{\operatorname{skew}\!\left( #1 \right)}
\newcommand{\symmop}{\operatorname{sym}}
\newcommand{\symm}[1]{\symmop\!\left( #1 \right)}
\newcommand{\symmsmall}[1]{\symmop( #1 )}
%\newcommand{\vecc}{\mathrm{vec}}
\newcommand{\veccc}[1]{\vecc\left({#1}\right)}
\newcommand{\sbdop}{\operatorname{symblockdiag}}
\newcommand{\sbd}[1]{\sbdop\!\left({#1}\right)}
\newcommand{\sbdsmall}[1]{\sbdop({#1})}
\newcommand{\RMSE}{\mathrm{RMSE}}
\newcommand{\Stdpm}{\St(d, p)^m}
\newcommand{\Proj}{\mathrm{Proj}}
\newcommand{\ProjH}{\mathrm{Proj}^h}
\newcommand{\ProjV}{\mathrm{Proj}^v}
\newcommand{\Exp}{\mathrm{Exp}}
\newcommand{\Log}{\mathrm{Log}}
\newcommand{\expect}{\mathbb{E}}
\newcommand{\expectt}[1]{\mathbb{E}\left\{{#1}\right\}}
%\newcommand{\R}{\mathrm{R}}
\newcommand{\Retr}{\mathrm{Retraction}}
\newcommand{\Trans}{\mathrm{Transport}}
\newcommand{\Transbis}[2]{\mathrm{Transport}_{{#1}}({#2})}
\newcommand{\polar}{\mathrm{polar}}
\newcommand{\qf}{\mathrm{qf}}
\newcommand{\Gr}{\mathrm{Gr}}
\newcommand{\St}{\mathrm{St}}
\newcommand{\GLr}{{\mathbb{R}^{r\times r}_*}}
\newcommand{\Precon}{\mathrm{Precon}}
\newcommand{\T}{\mathrm{T}}
\newcommand{\HH}{\mathrm{H}}
\newcommand{\VV}{\mathrm{V}}
\newcommand{\N}{\mathcal{N}}
\newcommand{\M}{\mathcal{M}}
\newcommand{\barM}{\overline{\mathcal{M}}}
\newcommand{\barnabla}{\overline{\nabla}}
\newcommand{\barf}{{\overline{f}}}
\newcommand{\barx}{{\overline{x}}}
\newcommand{\bary}{{\overline{y}}}
\newcommand{\barX}{{\overline{X}}}
\newcommand{\barY}{{\overline{Y}}}
\newcommand{\barg}{{\overline{g}}}
\newcommand{\barxi}{{\overline{\xi}}}
\newcommand{\bareta}{{\overline{\eta}}}
\newcommand{\p}{\mathcal{P}}
\newcommand{\SOn}{{\mathrm{SO}(n)}}
\newcommand{\On}{{\mathrm{O}(n)}}
\newcommand{\Od}{{\mathrm{O}(d)}}
\newcommand{\Op}{{\mathrm{O}(p)}}
\newcommand{\son}{{\mathfrak{so}(n)}}
\newcommand{\SOt}{{\mathrm{SO}(3)}}
\newcommand{\sot}{{\mathfrak{so}(3)}}
\newcommand{\SOtwo}{{\mathrm{SO}(2)}}
\newcommand{\sotwo}{{\mathfrak{so}(2)}}
\newcommand{\So}{{\mathbb{S}^{1}}}
\newcommand{\Stwo}{{\mathbb{S}^{2}}}
\newcommand{\Sp}{{\mathbb{S}^{2}}}
\newcommand{\Rtt}{{\mathbb{R}^{3\times 3}}}
\newcommand{\Snn}{{\mathbb{S}^{n\times n}}}
\newcommand{\Smm}{{\mathbb{S}^{m\times m}}}
\newcommand{\Smdmd}{{\mathbb{S}^{md\times md}}}
\newcommand{\Spp}{{\mathbb{S}^{p\times p}}}
\newcommand{\Sdd}{{\mathbb{S}^{d\times d}}}
\newcommand{\Srr}{{\mathbb{S}^{r\times r}}}
\newcommand{\Rnn}{{\mathbb{R}^{n\times n}}}
\newcommand{\Rnp}{{\mathbb{R}^{n\times p}}}
\newcommand{\Rdp}{{\mathbb{R}^{d\times p}}}
\newcommand{\Rmn}{{\mathbb{R}^{m\times n}}}
\newcommand{\Rnm}{\mathbb{R}^{n\times m}}
\newcommand{\Rnr}{\mathbb{R}^{n\times r}}
\newcommand{\Rrr}{\mathbb{R}^{r\times r}}
\newcommand{\Rpp}{\mathbb{R}^{p\times p}}
\newcommand{\Rmm}{{\mathbb{R}^{m\times m}}}
\newcommand{\Rmr}{\mathbb{R}^{m\times r}}
\newcommand{\Rrn}{{\mathbb{R}^{r\times n}}}
\newcommand{\Rd}{{\mathbb{R}^{d}}}
\newcommand{\Rp}{{\mathbb{R}^{p}}}
\newcommand{\Rk}{{\mathbb{R}^{k}}}
\newcommand{\Rdd}{{\mathbb{R}^{d\times d}}}
\newcommand{\RNN}{{\mathbb{R}^{N\times N}}}
\newcommand{\reals}{{\mathbb{R}}}
\newcommand{\complex}{{\mathbb{C}}}
\newcommand{\Rn}{{\mathbb{R}^n}}
\newcommand{\Rm}{{\mathbb{R}^m}}
\newcommand{\CN}{{\mathbb{C}^{\,N}}}
\newcommand{\CNN}{{\mathbb{C}^{\ N\times N}}}
\newcommand{\CnotR}{{\mathbb{C}\ \backslash\mathbb{R}}}
\newcommand{\grad}{\mathrm{grad}}
\newcommand{\Hess}{\mathrm{Hess}}
\newcommand{\vecc}{\mathrm{vec}}
\newcommand{\sign}{\mathrm{sign}}
\newcommand{\diag}{\mathrm{diag}}
\newcommand{\D}{\mathrm{D}}
\newcommand{\dt}{\mathrm{d}t}
\newcommand{\ds}{\mathrm{d}s}
\newcommand{\calA}{\mathcal{A}}
\newcommand{\calN}{\mathcal{N}}
\newcommand{\calM}{\mathcal{M}}
\newcommand{\calC}{\mathcal{C}}
\newcommand{\calE}{\mathcal{E}}
\newcommand{\calF}{\mathcal{F}}
\newcommand{\calL}{\mathcal{L}}
\newcommand{\calU}{\mathcal{U}}
\newcommand{\calH}{\mathcal{H}}
\newcommand{\calO}{\mathcal{O}}
\newcommand{\Id}{\operatorname{Id}} % need to be distinguishable from identity matrix
\newcommand{\rank}{\operatorname{rank}}
\newcommand{\nulll}{\operatorname{null}}
\newcommand{\col}{\operatorname{col}}
\newcommand{\Kmax}{K_{\mathrm{max}}}

\newcommand{\frobnormbig}[1]{\big\|{#1}\big\|_\mathrm{F}}
%\newcommand{\frobnorm}[1]{\left\|{#1}\right\|_\mathrm{F}}
\newcommand{\opnormsmall}[1]{\|{#1}\|_\mathrm{op}}
\newcommand{\frobnormsmall}[1]{\|{#1}\|_\mathrm{F}}
\newcommand{\smallfrobnorm}[1]{\|{#1}\|_\mathrm{F}}
\newcommand{\norm}[1]{\left\|{#1}\right\|}
\newcommand{\sqnorm}[1]{\left\|{#1}\right\|^2}
%\newcommand{\sqfrobnorm}[1]{\frobnorm{#1}^2}
\newcommand{\sqfrobnormsmall}[1]{\frobnormsmall{#1}^2}
\newcommand{\sqfrobnormbig}[1]{\frobnormbig{#1}^2}

\newcommand{\frobnorm}[2][F]{\left\|{#2}\right\|_\mathrm{#1}}
\newcommand{\sqfrobnorm}[2][F]{\frobnorm[#1]{#2}^2}

\newcommand{\dd}[1]{\frac{\mathrm{d}}{\mathrm{d}#1}}
\newcommand{\dmu}{\mathrm{d}\mu}
\newcommand{\pdd}[1]{\frac{\partial}{\partial #1}}
\newcommand{\TODO}[1]{{\color{red}{[#1]}}}
%\newcommand{\TODO}[1]{}
\newcommand{\dblquote}[1]{``#1''}
\newcommand{\floor}[1]{\lfloor #1 \rfloor}
\newcommand{\uniform}{\mathrm{Uni}}
\newcommand{\langevin}{\mathrm{Lang}}
\newcommand{\langout}{\mathrm{LangUni}}
\newcommand{\Zeq}{{Z_\mathrm{eq}}}
\newcommand{\feq}{{f_\mathrm{eq}}}
\newcommand{\XM}{\mathfrak{X}(\M)}
\newcommand{\FM}{\mathfrak{F}(\M)}
\newcommand{\length}{\operatorname{length}}
\newcommand{\bfR}{\mathbf{R}}
\newcommand{\bfxi}{{\bm{\xi}}}
\newcommand{\bfX}{{\mathbf{X}}}
\newcommand{\bfY}{{\mathbf{Y}}}
\newcommand{\bfeta}{{\bm{\eta}}}
\newcommand{\bfF}{\mathbf{F}}
\newcommand{\bfOmega}{\boldsymbol{\Omega}}
\newcommand{\bftheta}{{\bm{\theta}}}
\newcommand{\lambdamax}{\lambda_\mathrm{max}}
\newcommand{\lambdamin}{\lambda_\mathrm{min}}


\newcommand{\ddiag}{\mathrm{ddiag}}
\newcommand{\dist}{\mathrm{dist}}
\newcommand{\embdist}{\mathrm{embdist}}
\newcommand{\var}{\mathrm{var}}
\newcommand{\Cov}{\mathbf{C}}
\newcommand{\Covm}{C}
\newcommand{\FIM}{\mathbf{F}}
\newcommand{\FIMm}{F}
\newcommand{\Rmle}{{\hat\bfR_\mle}}
%\newcommand{\diag}{\mathbf{\mathrm{diag}}}
\newcommand{\Rmoptensor}{\mathbf{{R_\mathrm{m}}}}
\newcommand{\barRmoptensor}{\mathbf{{\bar{R}_\mathrm{m}}}}
\newcommand{\Rmop}{R_\mathrm{m}}
\newcommand{\barRmop}{\bar{R}_\mathrm{m}}
\newcommand{\Rcurv}{\mathcal{R}}
\newcommand{\barRcurv}{\bar{\mathcal{R}}}

\newcommand{\ith}{$i$th }
\newcommand{\jth}{$j$th }
\newcommand{\kth}{$k$th }
\newcommand{\ellth}{$\ell$th }

\newcommand{\expp}[1]{{\exp\!\big({#1}\big)}}


%\definecolor{darkgreen}{rgb}{0,.5,0}

\newtheorem{theorem}{Theorem} %[section]
\newtheorem{lemma}[theorem]{Lemma}
\newtheorem{proposition}[theorem]{Proposition}
\newtheorem{corollary}[theorem]{Corollary}
\newtheorem{conjecture}[theorem]{Conjecture}
\newtheorem{assumption}{Assumption} %[section]
\newtheorem{example}{Example} %[section]
\newtheorem{definition}{Definition} %[section]
\newtheorem{remark}{Remark}

%\newenvironment{proof}[1][Proof]{\begin{trivlist}
%\item[\hskip \labelsep {\bfseries #1}]}{\end{trivlist}}
%\newenvironment{definition}[1][Definition]{\begin{trivlist}
%\item[\hskip \labelsep {\bfseries #1}]}{\end{trivlist}}
%\newenvironment{example}[1][Example]{\begin{trivlist}
%\item[\hskip \labelsep {\bfseries #1}]}{\end{trivlist}}
%\newenvironment{remark}[1][Remark]{\begin{trivlist}
%\item[\hskip \labelsep {\bfseries #1}]}{\end{trivlist}}

%\newcommand{\qed}{\nobreak \ifvmode \relax \else
%      \ifdim\lastskip<1.5em \hskip-\lastskip
%      \hskip1.5em plus0em minus0.5em \fi \nobreak
%      \vrule height0.75em width0.5em depth0.25em\fi}
